\subsection{凸集的交。凸集的乘积}

命题 5. --- 仿射空间 E 的凸子集的任意一族之交是凸的。

本命题由 II, 第 7 页 的定义 1 立即得出。

命题 6. --- 设 \((\mathrm{E}_{\iota})_{\iota\in\mathbf{I}}\) 是向量空间的一个族，且对每个 \(\iota\in I\)，设 \(A_{\iota}\) 是 \(E_{\iota}\) 的一个非空子集。则集 \(A=\prod_{\iota\in I}A_{\iota}\) 在 \(E=\prod_{\iota\in I}E_{\iota}\) 中是凸的，当且仅当对所有 \(\iota\in I\)，集 \(A_{\iota}\) 在 \(E_{\iota}\) 中是凸的。

事实上，每个投影 \(\mathrm{pr}_{\iota}\) 都是线性映射，且我们有 \(A_{\iota} = \mathrm{pr}_{\iota}A\) 与 \(A = \bigcap_{\iota \in I} \mathrm{pr}_{\iota}^{-1}(A_{\iota})\)；本命题由上面的 prop. 2 与 prop. 5 得出。

推论. --- 在空间 \(\mathbf{R}^n\) 中，每个超平行体 (GT, VI, § 1.3) 都是凸子集。

因为它是矩形平行多面体在一个仿射线性映射下的像，而后者由 prop. 6 是凸的。

命题 7. --- 设 A 与 B 是向量空间 E 的两个凸子集。对任意实数 \(\alpha\)、\(\beta\)，集 \(\alpha A + \beta B\)（形如 \(\alpha x + \beta y\) 的点组成的集，其中 x 在 A 中变化，y 在 B 中变化）是凸的。

因为 \(\alpha A + \beta B\) 是凸子集 \(A \times B\)（\(E \times E\) 中的）在线性映射 \((x, y) \mapsto \alpha x + \beta y\)（\(E \times E\) 到 E 中的）下的像。

